\chapter{Membangun Tim AI: Strategi \& Kultur}
\label{chap:ai_team}

Mengadopsi AI bukan sekadar membeli lisensi ChatGPT. Ini adalah transformasi budaya yang memerlukan struktur tim yang tepat, pola pikir baru, dan kepemimpinan yang visioner. Bab ini membahas cara membangun "AI-Native Organization".

\section{Struktur Tim AI yang Efektif}

\subsection{1. Tim Terpusat (Center of Excellence)}
Semua talenta data (Data Scientist, MLE, Data Engineer) berada dalam satu departemen di bawah Chief Data Officer (CDO).
\begin{itemize}
    \item **Kelebihan:** Standarisasi kuat, berbagi pengetahuan (knowledge sharing) mudah, jenjang karir jelas.
    \item **Kekurangan:** Jauh dari masalah bisnis nyata, sering menjadi "pabrik eksperimen" yang tidak pernah masuk produksi.
\end{itemize}

\subsection{2. Tim Terdesentralisasi (Embedded)}
Setiap departemen (Marketing, Finance, Product) memiliki Data Scientist sendiri.
\begin{itemize}
    \item **Kelebihan:** Sangat dekat dengan masalah bisnis, iterasi cepat.
    \item **Kekurangan:** Silo data, duplikasi pekerjaan, tidak ada standar teknik yang konsisten.
\end{itemize}

\subsection{3. Model Hub-and-Spoke (Hybrid)}
Kombinasi terbaik. Ada tim inti (Hub) yang menyediakan infrastruktur dan standar, sementara Data Scientist (Spoke) ditempatkan di departemen bisnis tetapi melapor secara teknis ke Hub.
\begin{itemize}
    \item **Rekomendasi:** Gunakan model ini untuk perusahaan skala menengah-besar.
\end{itemize}

\section{Peran Kunci (The AI Squad)}

Untuk satu proyek AI end-to-end, Anda membutuhkan "squad" minimalis:
1.  **Product Manager (PM):** Mendefinisikan "Why". Mengapa kita butuh AI ini? Apa metrik suksesnya?
2.  **Machine Learning Engineer (MLE):** Membangun "How". Arsitektur model, deployment, skalabilitas.
3.  **Data Engineer:** Menyiapkan "Fuel". Pipa data (pipeline) yang bersih dan reliabel.
4.  **Software Engineer (Backend/Frontend):** Mengintegrasikan model ke aplikasi user.
5.  **Domain Expert (SME):** Dokter (untuk AI medis), Akuntan (untuk AI keuangan). Mereka yang memvalidasi output.

\section{Hiring: Menemukan Talenta Tepat}

\subsection{Red Flags saat Wawancara}
- **"Saya hanya mau pakai model terbaru (LLM)."** Hindari. Cari engineer yang mau membersihkan data kotor dan mencoba Regresi Logistik dulu.
- **"Saya tidak peduli bisnis, saya cuma mau riset."** Cocok untuk lab riset, tapi bahaya untuk startup/perusahaan produk.
- **Tidak bisa koding dasar.** Data Scientist yang tidak bisa menulis kode Python bersih akan menjadi beban bagi tim Engineering.

\subsection{Kualitas yang Dicari}
- **Curiosity:** AI berubah tiap minggu. Cari orang yang belajar mandiri.
- **Pragmatisme:** Memilih solusi sederhana yang bekerja daripada solusi kompleks yang keren.
- **Komunikasi:** Bisa menjelaskan konsep teknis ke orang awam.

\section{Membangun Budaya Data (Data Culture)}

\subsection{1. Demokratisasi Data}
Jangan kunci data di "menara gading". Berikan akses aman ke semua karyawan (misal: lewat Metabase atau Tableau). Biarkan mereka bertanya pada data.

\subsection{2. Toleransi Kegagalan (Fail Fast)}
Proyek AI bersifat probabilistik. Tidak semua eksperimen berhasil. Ciptakan lingkungan di mana gagal itu boleh, asalkan ada pelajaran yang diambil (post-mortem).

\subsection{3. Etika di Depan (Ethics First)}
Jangan biarkan etika menjadi renungan terakhir (afterthought). Diskusikan bias dan dampak sosial sejak hari pertama perencanaan proyek.

\section{Transformasi ke AI-Native}

Perusahaan AI-Native bukan hanya perusahaan yang "memakai AI", tapi yang "berpikir dengan AI".
- **Otomatisasi Default:** Setiap proses manual harus ditantang: "Bisakah AI membantu ini?"
- **Keputusan Berbasis Data:** Bukan berdasarkan "firasat bos" (HiPPO - Highest Paid Person's Opinion), tapi berdasarkan bukti data.

\section{Penutup: Peran Pemimpin}

Sebagai pemimpin, tugas Anda bukan menjadi ahli teknis, melainkan menjadi **Arsitek Lingkungan**. Tugas Anda adalah menyingkirkan hambatan birokrasi, menyediakan sumber daya komputasi (GPU/Cloud), dan memberikan visi yang jelas ke mana kapal ini berlayar.
